\addcontentsline{toc}{chapter}{KẾT LUẬN}
\chapter*{KẾT LUẬN}

\section*{Tóm tắt đóng góp}

Báo cáo nghiên cứu đã đề xuất và hiện thực một quy trình xây dựng mô hình ngôn ngữ nhỏ cho lĩnh vực pháp lý tiếng Việt, có khả năng giải quyết đồng thời ba tác vụ: suy luận tự nhiên (NLI), trắc nghiệm (MC) và suy luận tam đoạn luận. Đóng góp chính của nghiên cứu là sự kết hợp hai chiến lược sinh dữ liệu vốn được áp dụng riêng rẽ trong hai phương pháp dẫn đầu cuộc thi VLSP2025-LegalSML - sinh few-shot kiểu MinLegal~\cite{luan2025minlegal} và sinh theo khía cạnh (aspect-based) kiểu Bosch@AI~\cite{quang2025bosch} - trên cùng một baseline \texttt{qwen3-4b-legal-pretrain} và cùng điều kiện huấn luyện. Quy trình huấn luyện hai giai đoạn (QLoRA theo từng tác vụ rồi hợp nhất, sau đó tiếp tục tinh chỉnh adapter hợp nhất) được mô phỏng từ MinLegal và phân tích định lượng.

Một đóng góp thực dụng quan trọng là việc chứng minh toàn bộ quy trình có thể chạy được trên một GPU Tesla T4 16GB miễn phí (thực nghiệm trên Kaggle), nhờ kết hợp lượng tử hóa 4-bit NF4 theo cơ chế QLoRA~\cite{dettmers2023qlora}, gradient checkpointing và tích lũy gradient. Điều này hạ thấp đáng kể rào cản phần cứng so với các phương pháp gốc vốn cần cụm GPU A30/A100. Ngoài ra, việc sử dụng Claude Haiku 4.5 thay cho GPT-4o trong cả vai trò sinh dữ liệu lẫn mô hình giám khảo giúp giảm chi phí khiến nghiên cứu dễ tái lập hơn cho cộng đồng sinh viên và nhà nghiên cứu hạn chế tài nguyên.

Về mặt thực nghiệm, mô hình đề xuất đạt MC=86,99\%, NLI=89,33\%, Syllogism=56,70\% với điểm trung bình 77,67\% trên tập đánh giá (440 mẫu gốc). Tuy điểm trung bình còn thấp hơn hai phương pháp gốc, mô hình đề xuất vượt Bosch@AI ở riêng tác vụ tam đoạn luận (56,70\% so với 53,58\%) dù dùng ít dữ liệu hơn nhiều - minh chứng cho hiệu quả của chiến lược sinh dữ liệu theo khía cạnh (aspect-based) đối với năng lực lập luận. Phân tích định tính đóng góp của từng chiến lược sinh dữ liệu cũng cung cấp hiểu biết về việc chiến lược nào phù hợp với từng loại tác vụ.

\section*{Hạn chế}

Nghiên cứu còn một số hạn chế. Thứ nhất, quy mô dữ liệu tổng hợp nhỏ hơn đáng kể so với hai phương pháp gốc (vài nghìn so với hàng trăm nghìn mẫu), do ràng buộc về chi phí API và tài nguyên GPU miễn phí. Thứ hai, việc chọn baseline \texttt{qwen3-4b-legal-pretrain} là một đánh đổi có chủ ý: MinLegal báo cáo phiên bản Qwen3-4B-Base có thể cho kết quả nhỉnh hơn nhưng nghiên cứu ưu tiên kế thừa giai đoạn domain pre-training để tiết kiệm chi phí. Thứ ba, độ tin cậy của mô hình giám khảo Claude Haiku 4.5 đối với tác vụ tam đoạn luận chưa được kiểm chứng đầy đủ so với giám khảo con người, có thể ảnh hưởng đến điểm số của tác vụ này. Cuối cùng, tam đoạn luận vẫn là điểm yếu chung, phản ánh thách thức cố hữu của lập luận pháp lý đa bước.

\section*{Hướng phát triển}

Trong tương lai, nghiên cứu có thể được mở rộng theo nhiều hướng. Về dữ liệu, có thể mở rộng kho ngữ liệu sinh và kết hợp truy xuất điều luật (retrieval-augmented generation) để cung cấp ngữ cảnh chính xác cho mô hình, giảm ảo giác. Về đánh giá, có thể xây dựng một mô hình giám khảo đáng tin cậy hơn hoặc đối chiếu với chú thích của chuyên gia pháp lý cho tác vụ tam đoạn luận. Về kỹ thuật tinh chỉnh, có thể thử nghiệm các biến thể tiên tiến của LoRA~\cite{hu2022lora} như DoRA hoặc AdaLoRA để cải thiện hiệu quả tham số. Những hướng này hứa hẹn tiếp tục thu hẹp khoảng cách hiệu năng với các phương pháp dẫn đầu trong khi vẫn duy trì chi phí thấp.
